Generating synthetic financial time series that faithfully preserve the statistical properties of real market data is a central challenge in quantitative finance.
High-fidelity synthetic data are essential for stress testing risk models against plausible but unobserved market scenarios, validating portfolio optimization algorithms, and augmenting limited training sets for machine learning pipelines \citep{assefa2020generating, jordon2022synthetic}.
Empirical equity excess growth rates exhibit three well-documented statistical regularities---heavy-tailed (leptokurtic) distributions, negligible linear autocorrelation in raw returns, and persistent volatility clustering---that any generative framework must simultaneously reproduce to be considered statistically faithful \citep{cont2001empirical, mandelbrot1963variation}.

\citet{alswaidan2026hybrid} addressed this challenge with a discrete hidden Markov model (HMM) that partitioned continuous excess growth rates into Laplace quantile-defined states and estimated transitions by direct frequency counting.
A Poisson jump-duration mechanism forced the model to dwell in high-volatility tail states for empirically realistic durations, partially reproducing volatility clustering.
That framework achieved high distributional pass rates and scaled to a 424-asset pipeline, but the discretization step introduced quantization error: continuous observations were binned into categorical states, and the emission model sampled from within-bin Student-$t$ distributions rather than learning the emission parameters jointly with the transition structure.

In this study, we extend that framework to \emph{continuous Gaussian emissions} learned via the Baum-Welch (Expectation-Maximization) algorithm \citep{baum1970maximization, rabiner1986introduction}.
Each hidden state emits from a Normal distribution $\mathcal{N}(\mu_k, \sigma_k)$ whose parameters are estimated iteratively alongside the transition matrix, eliminating the need for quantile binning and enabling the model to capture fine-grained regime dynamics directly from the data.
The Poisson jump-duration mechanism is preserved, operating on the continuous model's state space to enforce tail-state persistence.

The key contributions of this work are:
\begin{enumerate}
    \item \textbf{Baum-Welch training} of emission means, variances, and transition probabilities directly from continuous observations, with a quantile-based initialization strategy that avoids degenerate local optima.
    \item \textbf{State resolution analysis} across $K \in \{3, 6, 9, 11, 13\}$ hidden states, demonstrating that distributional fidelity improves monotonically with $K$ while temporal structure (ACF-MAE) remains stable.
    \item \textbf{Comprehensive validation} using five complementary metrics---Kolmogorov-Smirnov (KS) pass rates, excess kurtosis matching, autocorrelation mean absolute error (ACF-MAE), Wasserstein-1 distance, and Hellinger distance---on both in-sample (2014--2024) and out-of-sample (2025) data.
    \item \textbf{Cross-asset generalization} to NVDA, JNJ, and JPM, confirming that the framework adapts to diverse equity risk profiles without modification.
    \item \textbf{Extension to volatility indices}: application of the CHMM framework to CBOE VIX log returns, an asset class with distinct statistical properties (positive skewness, mean reversion, extreme kurtosis), demonstrating that the framework generalizes beyond equity returns to volatility dynamics. The VIX analysis further clarifies the role of the jump mechanism: at small $K$, the base model already captures volatility clustering, but as $K$ increases the learned transition matrix produces insufficient tail-state persistence, making the Poisson jump extension essential.
    \item \textbf{Pathway to implied volatility modeling}: the regime-dependent emission variances learned by the CHMM provide a natural parameterization for the mean-reversion target $\theta(S_t)$ in a modified Heston stochastic volatility model, enabling closed-form implied volatility generation for option pricing without circular dependence on market-observed IV surfaces.
\end{enumerate}

The continuous emission structure provides a natural foundation for stochastic volatility extensions that the discrete quantile-binning approach cannot support.
In particular, the VIX application demonstrates that the CHMM framework can directly model the volatility process itself---not merely generate synthetic returns---opening a pathway to regime-aware option pricing via integration with the Heston model \citep{heston1993closed}.
